% mode: ini-etex
% Aimed at scan_file_name, more_name, end_name and start_input: every way of writing a name to \input, with the
% name that comes out shown by \jobname, \inputlineno and the parenthesis in the log.  The empty file, a file with no
% last newline, CRLF and trailing blanks go through input_ln.  The braces go through scan_file_name_braced.
\input prelude
\def\x{files/x}\def\ext{.inc}\def\space{ }
\message{<\jobname>}
% Plain forms, ended by a space, by the end of the line, and by something that is not a character.
\input files/x.inc
\input files/x.inc \message{<space ended it>}
\input files/x.inc\relax\message{<relax ended it>}
\input files/y.inc
% Spaces before the name are skipped, inside are not.
\input    files/x.inc   \message{<many spaces>}
\input
   files/x.inc
\message{<name on the next line>}
% Macros: the name is expanded as it is scanned.
\input\x.inc
\input\x\ext
\input\x\ext\relax
\expandafter\input\x.inc
% Quotes: anywhere in the name, around a part of it, empty, and running to the end of the line.
\input "files/x.inc"
\input files/"x.inc"
\input fi"le"s/x.inc
\input files/x"".inc
\input "files/with space.inc"
\input "files/with space".inc
\input files/"with space".inc
\input files/with" "space.inc
\input "files/with space.inc
\message{<quote ran to the end of the line>}
\def\q{"}
\input \q files/with space.inc\q
\message{<quotes from a macro>}
% Odd contents.
\input files/empty.inc
\input files/nonl.inc
\input files/crlf.inc
\input files/spaces.inc
\input files/endinput.inc
% The extension .tex is tried first, then the name as it stands: dot.tex here.
\input files/dot
\input files/dot.tex
% An area with dots, and a ./
\input ./files/x.inc
\input files/../files/x.inc
\input files/./x.inc
% An \input in the middle of a name is not expanded: \relax goes in front, and it ends the name.
\input files/x.inc\input files/x.inc
\message{<inner input>}
% \expandafter and \csname during a name; a new control sequence must not disturb the string in progress.
\setbox0\hbox{\input files/dot\csname brandnewname\endcsname .inc}
\show\brandnewname
\input files/x.inc\csname relax\endcsname
% Braces: spaces and quotes are not special, or quotes are removed all the same.
\input{files/x.inc}
\input {files/x.inc}\message{<braces and a space before>}
\input
{files/x.inc}
\input{\x\ext}
\input{\x.inc}
\input{files/with space.inc}
\input{"files/with space.inc"}
\input{files/"with space".inc}
\expandafter\input\expandafter{\x.inc}
\input{files/x\noexpand.inc}
% In \edef, in \csname: file name scanning from a token list.
\edef\y{\noexpand\input\space files/x.inc\space}\show\y \y
\def\z{\input files/x.inc}\z\message{<at the end of a macro>}
\def\z{\input files/x.inc }\z\message{<space at the end of a macro>}
\def\z#1{\input #1}\z{files/x.inc}\message{<from an argument>}
\expandafter\ifx\csname undefined\endcsname\relax \input files/x.inc \fi
\ifnum0=0 \input files/x.inc\fi
\message{<in a conditional>}
\begingroup \input files/x.inc\endgroup
% \scantokens opens a pseudo-file with no name; \input inside it.
\scantokens{\input files/x.inc}
\scantokens{\message{<in scantokens>}\input files/x.inc \message{<after it>}}
% The input stack, as show_context prints it, with names and lines.
\input files/showcontext.inc
\message{<inputlineno \the\inputlineno>}
\end
